El artículo \textit{Offensive Lineup Analysis in Basketball with Clustering Players Based on Shooting Style and Offensive Role}, desarrollado por Yamada y Fujii (2024), se centra en analizar la compatibilidad entre jugadores de baloncesto a partir de sus estilos de tiro y las funciones ofensivas que desempeñan dentro de la cancha. Los autores plantean que el rendimiento ofensivo de un equipo no depende únicamente de contar con buenos anotadores, sino también de identificar qué tipos de jugadores se complementan mejor al formar una alineación.

Para realizar este análisis utilizan dos enfoques diferentes. El primero consiste en clasificar a los jugadores según sus estilos de tiro, mientras que el segundo busca identificar sus roles ofensivos a partir de las acciones que realizan durante los partidos. Posteriormente, utilizan estas clasificaciones para evaluar la eficiencia de las combinaciones de jugadores.

En el análisis de estilos de tiro utilizaron datos de seguimiento espacial (\textit{tracking}) de 630 partidos de la temporada 2015--2016 de la NBA, considerando un total de 41,160 lanzamientos. De cada tiro extrajeron 17 características relacionadas con la posición del jugador, la distancia al aro, el desplazamiento y el tiempo que mantiene el balón antes de realizar el lanzamiento.

Para reducir la cantidad de variables utilizaron el análisis de componentes principales (PCA), pasando de 17 características a 9 dimensiones. Después calcularon la distancia de Wasserstein para comparar las distribuciones de tiros de los jugadores y aplicaron agrupamiento jerárquico mediante el método de Ward. Como resultado de este procedimiento identificaron 13 grupos de jugadores con diferentes estilos de tiro.

Para el segundo enfoque utilizaron estadísticas ofensivas de las temporadas 2015--2016 hasta 2022--2023 de la NBA. En este caso analizaron diferentes tipos de jugadas, como:

\begin{itemize}
    \item Jugadas de aislamiento (\textit{Isolation}).
    \item Bloqueos directos (\textit{Pick-and-roll}).
    \item Jugadas de transición (\textit{Transition}).
    \item Lanzamientos tras recepción (\textit{Spot-up}).
    \item Movimientos sin balón (\textit{Off-screen}).
    \item Jugadas cercanas al aro (\textit{Post-up}).
\end{itemize}

Inicialmente aplicaron el algoritmo K-means para agrupar a los jugadores según sus características ofensivas. Sin embargo, identificaron que algunos jugadores podían desempeñar diferentes funciones durante un partido y que asignarlos a una sola categoría no representaba completamente su estilo de juego.

Para resolver este problema utilizaron el algoritmo \textit{Fuzzy C-means}, que permite que un jugador pertenezca a diferentes grupos con distintos grados de pertenencia. De esta manera obtuvieron 10 grupos de roles ofensivos, considerando que un mismo jugador puede combinar características de diferentes estilos de juego.

Una vez construidos los perfiles, los autores analizaron la eficiencia de las alineaciones utilizando el indicador \textit{Offensive Rating} (OFFRTG), que representa los puntos anotados por cada 100 posesiones. Para estudiar las combinaciones de cinco jugadores utilizaron modelos de aprendizaje automático, mientras que para evaluar la compatibilidad entre parejas de estilos ofensivos emplearon modelos bayesianos.

Los resultados muestran que ciertas combinaciones de estilos de juego tienen una relación positiva con la eficiencia ofensiva del equipo. Por ejemplo, se identificaron combinaciones favorables entre jugadores que pueden generar oportunidades de anotación y aquellos que se especializan en lanzamientos exteriores o en facilitar la circulación del balón. Por otra parte, algunas combinaciones de jugadores con funciones ofensivas similares presentaron efectos menos favorables.

Una de las principales ventajas de este estudio es que no utiliza únicamente las posiciones tradicionales de los jugadores, como base, escolta, alero o pívot, sino que construye perfiles a partir de sus características y acciones ofensivas. Esto permite identificar diferencias entre jugadores que ocupan una misma posición, pero que desempeñan funciones distintas durante el partido.

Otra ventaja es el uso de \textit{Fuzzy C-means}, ya que permite representar jugadores con características híbridas, evitando que cada uno tenga que pertenecer obligatoriamente a una sola categoría. Además, el uso de modelos bayesianos permite cuantificar los efectos de las combinaciones de estilos ofensivos y comparar cuáles presentan relaciones más favorables con el rendimiento del equipo.

Sin embargo, la principal limitación del artículo es que se centra únicamente en la eficiencia ofensiva, dejando de lado el impacto defensivo de los jugadores y de las alineaciones. Esto es importante debido a que una combinación puede producir una gran cantidad de puntos, pero también permitir que el equipo contrario anote con facilidad.

Otra limitación es que el análisis de estilos de tiro utiliza datos de seguimiento de una sola temporada, mientras que el análisis de roles ofensivos considera varias temporadas. Además, las combinaciones se evalúan principalmente a partir de los estilos de los jugadores, sin estudiar de manera directa cómo estos se comportan frente a las características particulares de una alineación rival.

También se debe considerar que los autores analizan la compatibilidad entre grupos de estilos de juego, por lo que los resultados no necesariamente representan el rendimiento específico de una pareja de jugadores. Dos jugadores pueden pertenecer a categorías similares y presentar resultados diferentes debido a sus capacidades individuales.
 \cite{yamada2024offensivelineupanalysisbasketball}.
